\documentclass[../../main.tex]{subfiles}
\begin{document}
\section{Lie 子代数与 Lie 子群}\label{section:ch19-lie-subalgebras}

叶状结构在 Lie 群理论中有深刻的应用。这里我们给出两个这样的应用；在接下来两章中还会看到更多。二者都依赖于 Lie 子代数与分布之间如下简单关系。若 Lie 群 $G$ 上的分布在每个左平移下都不变，则称它是\emph{左不变的}（回忆这意味着对每个 $g\in G$ 都有 $\mathrm{d}(L_g)(D)=D$）。

\begin{lemma}\label{lemma:ch19-子代数给出左不变分布}
设 $G$ 是一个 Lie 群。若 $\mathfrak h$ 是 $\operatorname{Lie}(G)$ 的一个 Lie 子代数，则子集
\begin{align}
D=\bigcup_{g\in G}D_g\subseteq TG,\qquad D_g=\{X_g\colon X\in\mathfrak h\}\subseteq T_gG
\end{align}
是 $G$ 上一个左不变的对合分布。
\end{lemma}

\begin{proof}
每个 $X\in\mathfrak h$ 都是 $G$ 上的一个左不变向量场。于是对任意 $g,g'\in G$，微分 $\mathrm{d}(L_{g'g^{-1}})$ 限制为从 $D_g$ 到 $D_{g'}$ 的同构。由此 $D_g$ 对每个 $g$ 维数相同，且 $D$ 是左不变的。$\mathfrak h$ 的任意一组基 $(X_1,\dots,X_k)$ 都是 $D$ 的一组整体光滑标架，故 $D$ 光滑。此外，由于对所有 $i,j\in\{1,\dots,k\}$ 都有 $[X_i,X_j]\in\mathfrak h$，由引理\cref{lemma:ch19-局部标架对合判别} 可知 $D$ 对合。
\end{proof}

\begin{theorem}[Lie 子群是弱嵌入的]\label{theorem:ch19-子群弱嵌入}
每个 Lie 子群都是一个对合分布的积分流形，因而是一个弱嵌入子流形。
\end{theorem}

\begin{proof}
设 $G$ 是一个 Lie 群，$H\subseteq G$ 是一个 Lie 子群。定理\cref{theorem:ch08-46} 表明 $H$ 的 Lie 代数典范同构于 Lie 子代数 $\mathfrak h=\iota_*\bigl(\operatorname{Lie}(H)\bigr)\subseteq\operatorname{Lie}(G)$，其中 $\iota\colon H\hookrightarrow G$ 是包含映射。令 $D\subseteq TG$ 为由 $\mathfrak h$ 按引理\cref{lemma:ch19-子代数给出左不变分布} 确定的对合分布。由定义可知，在每一点 $h\in H$ 处，切空间 $T_hH$ 等于 $D_h$，故 $H$ 是 $D$ 的一个积分流形。于是由定理\cref{theorem:ch19-弱嵌入}，$H$ 在 $G$ 中弱嵌入。
\end{proof}

\begin{theorem}[与 Lie 子代数相伴的 Lie 子群]\label{theorem:ch19-子代数对应子群}
设 $G$ 是一个 Lie 群，$\mathfrak g$ 是它的 Lie 代数。若 $\mathfrak h$ 是 $\mathfrak g$ 的任意 Lie 子代数，则存在 $G$ 的唯一一个连通 Lie 子群，其 Lie 代数为 $\mathfrak h$。
\end{theorem}

\begin{proof}
设 $\mathfrak h$ 是 $\mathfrak g$ 的一个 Lie 子代数。令 $D\subseteq TG$ 为由上式定义的对合分布。令 $\mathscr H$ 表示 $D$ 所确定的叶状结构；对任意 $g\in G$，令 $H_g$ 为 $\mathscr H$ 中含 $g$ 的叶（见图\cref{figure:ch19-subgroup}）。由于 $D$ 左不变，由命题\cref{proposition:ch19-不变分布与叶状} 可知每个左平移都把叶映到叶：对任意 $g,g'\in G$，有 $L_g(H_{g'})=H_{gg'}$。

定义 $H=H_e$，即含单位元的叶。我们将证明 $H$ 就是所求的 Lie 子群。

首先，为看出 $H$ 是一个子群，注意对任意 $h,h'\in H$，
\begin{align*}
hh'=L_h(h')\in L_h(H)=L_h(H_e)=H_h=H.
\end{align*}
于是乘法 $\psi\colon H\times H\to G$ 立即光滑。由于 $H$ 是一个对合分布的积分流形，定理\cref{theorem:ch19-弱嵌入} 表明它弱嵌入，故 $\psi$ 作为到 $H$ 的映射也是光滑的。$H$ 是 $\mathscr H$ 的一叶这一事实意味着 $H$ 的 Lie 代数为 $\mathfrak h$，因为 $H$ 在单位元处的切空间是 $D_e=\{X_e\colon X\in\mathfrak h\}$。为看出 $H$ 是唯一的 Lie 代数为 $\mathfrak h$ 的连通子群，设 $\widetilde H$ 是任意另一个 Lie 代数相同的连通子群。容易看出任意这样的 Lie 子群都是 $D$ 的积分流形，故由 $H=H_e$ 的极大性，必有 $\widetilde H\subseteq H$。另一方面，若 $U$ 是含单位元的 $D$ 的一个平坦坐标图的定义域，则由命题\cref{proposition:ch19-局部结构}，$\widetilde H\cap U$ 是切片开子集的并。由于含 $e$ 的切片是 $H$ 的一个开子集，这蕴含 $\widetilde H$ 含 $H$ 中单位元的一个邻域。由于单位元的任意邻域生成 $H$（命题\cref{proposition:ch07-14}），这蕴含 $\widetilde H=H$。
\end{proof}

\begin{figure}[ht]
\centering
\includegraphics[width=0.8\textwidth]{fig-19-9.png}
\caption{寻找 Lie 代数为 $\mathfrak{h}$ 的子群}\label{figure:ch19-subgroup}
\end{figure}
\end{document}
